\begin{abstract}
The ProveIt iPRA repository now contains a checked object calculus, a substantial library of
ordinary derivations, and more than one hundred mathematical proof packets.  The packets repeatedly
exhibit the same mismatch: finite search, exact arithmetic, trace checking, and certificate
verification are easy to describe at the host level, but the final claim is often not a term of the
canonical judgment \(\Prf\,\PRAAx\,\Gamma\,\varphi\).  Several prior frameworks solve the semantic
part of the problem, yet leave a generic last-mile compiler as a blueprint.

This report proposes \TAPIR-iPRA (\emph{Typed Arithmetic Proof Intermediate Representation}).  Its
central idea is deliberately modest: do not internalize a universal syntax interpreter inside iPRA.
Instead, specialize every mathematical construction at Lean elaboration time to concrete \(\PR\),
\(\PTm\), and \(\Fm\) expressions, and represent the resulting argument by an indexed proof plan
whose type already states its object-theory endpoint.  A total interpreter from plans to
\(\Prf\,\PRAAx\) is then a structural recursion using only checked constructors and checked library
lemmas.  An untrusted host generator may build plans, but it cannot manufacture a theorem because
Lean must accept the generated term.

The second ingredient is a typed sequential schedule.  A new node is allowed to use earlier outputs
and a fixed list of residual producer hypotheses.  Simultaneous cut discharges all earlier outputs,
so induction on the schedule produces a proof of every node from the residuals alone.  Consequently,
a finite contract graph compiles to an ordinary sequent
\[
  B_1,\ldots,B_r,\Gamma\;\vdash_{\PRAAx}\;\Phi,
\]
with no hidden oracle.  Supplying object proofs of the \(B_i\) closes the theorem; otherwise the
conditional endpoint precisely records what mathematics is still missing.

The report integrates certificate escape analysis, bounded least-witness extraction, compositional
bounded reflection, clocked traces, quantitative names, stable observations, and height-bounded
enumeration.  It resolves several difficulties that older packets had left open by reusing current
checked modules: \code{Prf.subst}, \code{Prf.weaken}, \code{Prf.trans}, quantifier-free induction,
least bounded search, bounded universal and existential reflection, pairing, and measure-decreasing
termination.  It does not claim ineffective height bounds, convergence moduli, separation constants,
or uniform generator algorithms.  Those remain explicit producer boundaries rather than being
smuggled into a purported compiler.

The archive contains the complete \LaTeX{} source and PDF, Lean-oriented modules and a
formalization plan, an executable reference implementation of graph slicing, schedule validation,
escape analysis, finite strategy synthesis, demand closure, and traces, together with tests and
machine-readable metadata.
\end{abstract}

\section*{Main contribution at a glance}
\addcontentsline{toc}{section}{Main contribution at a glance}

\begin{center}
\resizebox{0.97\textwidth}{!}{%
\begin{tikzpicture}[
  node distance=8mm and 12mm,
  box/.style={draw=tapirblue,rounded corners,fill=tapirgray,align=center,inner sep=6pt},
  arr/.style={-{Stealth[length=2.2mm]},thick,tapirblue}]
\node[box] (math) {Theorem-specific\\finite mathematics};
\node[box,right=of math] (specialize) {Host specialization\\to fixed terms and formulas};
\node[box,right=of specialize] (plan) {Typed \(\Plan\)\\and schedule};
\node[box,right=of plan] (prf) {Literal\\\(\Prf\,\PRAAx\)};
\draw[arr] (math) -- (specialize);
\draw[arr] (specialize) -- (plan);
\draw[arr] (plan) -- node[above,align=center]{structural\\interpreter} (prf);
\node[box,below=of plan] (res) {Explicit residuals\\height bounds, moduli, gaps};
\draw[arr] (res) -- (plan);
\end{tikzpicture}%
}
\end{center}

The framework is not a stronger theory.  It is a disciplined way to obtain ordinary iPRA
proofs from the finite constructions that iPRA can actually express, while making every stronger
input visible in the context.
